\otazka{15. Vymezení bezpečnostních studií}{Pojmy bezpečnost, hrozba a riziko a jejich použití v oblasti kyberbezpečnosti. Vnitřní a vnější bezpečnost a sektory z pohledu kyberbezpečnosti. Vymezení bezpečnostní politiky a její analýza se zaměřením na politiku v oblasti kyberbezpečnosti.}{BSSb1101}

\subsection*{15. Vymezení bezpečnostních studií -- stručněji}
Tuhle otázku je nejlepší říct jako základní pojmový rámec. Nejdřív vysvětlím, co bezpečnostní studia zkoumají, potom pojmy bezpečnost, hrozba a riziko, potom vnitřní/vnější bezpečnost a nakonec bezpečnostní politiku na příkladu kyberbezpečnosti.

\textbf{1. Bezpečnostní studia}
\begin{itemize}
  \item Bezpečnostní studia zkoumají, jak aktéři rozpoznávají hrozby, jak chrání své hodnoty a jakými politickými, právními, vojenskými, ekonomickými, technologickými a společenskými nástroji bezpečnost zajišťují.
  \item Nejde jen o válku. Moderní bezpečnost zahrnuje stát, společnost, jednotlivce, kritickou infrastrukturu, ekonomiku, informace, technologie a kyberprostor.
\end{itemize}

\textbf{2. Bezpečnost, hrozba a riziko}
\begin{itemize}
  \item Bezpečnost je stav nebo proces, kdy jsou chráněny důležité hodnoty a funkce před ohrožením. Může být objektivní i subjektivní; člověk nebo stát se může cítit bezpečněji či méně bezpečně, než odpovídá realitě.
  \item Hrozba je potenciální zdroj škody: aktér, jev nebo událost, která může narušit chráněnou hodnotu. V kyberbezpečnosti to může být ransomware skupina, státní APT, phishing, insider nebo zranitelnost v dodavateli.
  \item Riziko je pravděpodobnost a dopad toho, že hrozba využije zranitelnost a způsobí škodu. Prakticky: hrozba je ransomware, zranitelnost jsou neaktualizované servery, riziko je šance a dopad zašifrování nemocnice.
\end{itemize}

\textbf{3. Vnitřní, vnější a sektorová bezpečnost}
\begin{itemize}
  \item Vnější bezpečnost řeší hrozby z mezinárodního prostředí, například cizí státy, válku, špionáž, sankční nátlak nebo hybridní působení.
  \item Vnitřní bezpečnost řeší pořádek a fungování státu uvnitř: kriminalitu, ochranu obyvatelstva, veřejnou správu, kritické služby a krizové řízení.
  \item V kyberprostoru se hranice stírá. Útok může přijít ze zahraničí, vést přes infrastrukturu třetích zemí, zasáhnout českou nemocnici a současně mít trestní, politickou i bezpečnostní rovinu.
  \item Sektory bezpečnosti lze říct jednoduše: vojenský, politický, ekonomický, societální, environmentální a kybernetický/informační rozměr. Kyberbezpečnost je průřezová, protože může zasáhnout všechny ostatní sektory.
\end{itemize}

\textbf{4. Bezpečnostní politika}
\begin{itemize}
  \item Bezpečnostní politika je vrcholný program státu nebo jiného aktéra, jak chce chránit své zájmy. Vychází z definice zájmů, analýzy bezpečnostního prostředí, klasifikace hrozeb a rizik a mezinárodních závazků.
  \item Deklaruje, jakými prostředky bude aktér své zájmy hájit. U kyberbezpečnosti to znamená strategie, zákony, instituce, technická opatření, vzdělávání, mezinárodní spolupráci a reakce na incidenty.
  \item Analýza bezpečnostní politiky se ptá: jaký problém stát vidí, jaký má cíl, jaké nástroje volí, kdo je odpovědný, jaké jsou náklady a jak poznáme, že politika funguje.
\end{itemize}

\subsection*{Kostra odpovědi}
\begin{enumerate}
  \item Vymezit bezpečnostní studia a vztah ke strategickým studiím.
  \item Vysvětlit pojmy bezpečnost, kybernetická bezpečnost, hrozba, zranitelnost, riziko a incident.
  \item Ukázat rozdíl mezi vnitřní a vnější bezpečností a proč se v kyberprostoru prolínají.
  \item Popsat pět sektorů Kodaňské školy a stručně vysvětlit sekuritizaci a bezpečnostní dilema.
  \item Vymezit bezpečnostní politiku a praktickou šablonu její analýzy v kyberbezpečnosti.
\end{enumerate}

\subsection*{Bezpečnostní studia}
Bezpečnostní studia zkoumají, co je ohroženo, kdo nebo co to ohrožuje a jakými prostředky lze bezpečnost chránit nebo obnovit. Většinou se chápou jako podobor mezinárodních vztahů. Vznikala přibližně v polovině 20. století, hlavně v USA a západní Evropě, a během studené války se soustředila zejména na válku, vojenskou bezpečnost, jaderné odstrašení a vztahy mezi státy.

Po skončení studené války se pojetí bezpečnosti rozšířilo. Vedle vojenských hrozeb se začaly řešit také hrozby ekonomické, politické, společenské, ekologické, kriminální, teroristické a kybernetické. Po 11. září 2001 výrazně posílila studia vnitřní bezpečnosti, protože se ukázalo, že bezpečnost státu neohrožuje jen armáda jiného státu, ale i terorismus, extremismus, kriminalita, radikalizace nebo selhání kritických služeb.

Hranice bezpečnostních studií nejsou úplně ostré. Obor využívá politologii, právo, ekonomii, sociologii, psychologii, historii i technické vědy. Kyberbezpečnost je dobrý příklad: útok na nemocnici je současně technický incident, právní otázka, problém krizového řízení, ekonomická škoda a ohrožení důvěry veřejnosti.

Strategická studia jsou užší. Zaměřují se hlavně na vojenskou bezpečnost, válku, strategii, ozbrojené síly a odstrašení. Bezpečnostní studia jsou širší: zahrnují i nevojenské a vnitřní hrozby, například kyberkriminalitu, ochranu kritické infrastruktury nebo bezpečnost veřejných služeb.

\subsection*{Bezpečnost a kybernetická bezpečnost}
Bezpečnost je stav, kdy jsou hrozby pro určitý objekt a jeho zájmy omezeny na přijatelnou úroveň a objekt je připraven zbývající hrozby zvládat. Vždy se proto ptáme: bezpečnost koho nebo čeho? Chráněným objektem může být stát, člověk, firma, nemocnice, informační systém, digitální služba nebo data.

Absolutní bezpečnost neexistuje. Bezpečnost znamená, že hrozby nejsou nepřijatelné a že objekt umí přežít, fungovat a dále se rozvíjet. Nemocnice například není bezpečná jen tehdy, když na ni nikdo nikdy nezaútočí. Bezpečná je tehdy, když má přiměřenou prevenci, umí incident včas zjistit, omezit jeho dopad a obnovit zdravotní péči.

Krátce je dobré rozlišit dvě roviny. \textbf{Objektivní bezpečnost} je skutečná úroveň hrozeb, zranitelností a ochrany. \textbf{Subjektivní bezpečnost} je pocit bezpečí nebo ohrožení. Firma může mít objektivně slabé zabezpečení, ale vedení má pocit, že je vše v pořádku, protože se zatím nic nestalo. Naopak mediálně viditelný incident může vyvolat silný pocit ohrožení i tam, kde bylo konkrétní riziko rychle omezeno.

Rozdíl mezi \emph{security} a \emph{safety} je praktický. \emph{Security} míří hlavně na ochranu před úmyslným jednáním, například hackingem, špionáží nebo sabotáží. \emph{Safety} míří na nehody, poruchy, lidské chyby a přírodní události. V kyberbezpečnosti se prolínají: ransomware je security problém, ale výpadek nemocničního systému může ohrozit safety pacientů.

Kybernetická bezpečnost je soubor právních, organizačních, technických a vzdělávacích opatření, která chrání systémy, sítě, aplikace, data a uživatele a umožňují zvládnout incident a obnovit provoz. \textbf{Právní} část znamená pravidla, povinnosti a odpovědnost. \textbf{Organizační} část znamená role, procesy, řízení rizik a kontinuitu. \textbf{Technická} část zahrnuje řízení přístupů, segmentaci, šifrování, monitoring a zálohy. \textbf{Vzdělávací} část znamená školení uživatelů i odborníků.

Prakticky se chráněné vlastnosti často vysvětlují přes CIA triádu: \textbf{důvěrnost} znamená, že se k datům dostane jen oprávněný, \textbf{integrita} znamená, že data a systémy nejsou neoprávněně změněny, a \textbf{dostupnost} znamená, že služba funguje, když ji potřebujeme. U kritických služeb se k tomu přidává odolnost a schopnost obnovy.

\subsection*{Hrozba, zranitelnost, riziko a incident}
Hrozba je jev, událost nebo aktér, který může nebo chce poškodit určitou hodnotu. V kyberbezpečnosti může být hrozbou ransomwarová skupina, státem podporovaná špionáž, DDoS útok, insider, lidská chyba, požár serverovny nebo technická porucha. Úmyslná hrozba je například ransomware, špionáž nebo DDoS. Neúmyslná hrozba je například chyba administrátora, požár nebo výpadek napájení.

Nezáplatovaný systém není hrozba, ale \textbf{zranitelnost}. Zranitelnost je slabina, přes kterou může hrozba způsobit škodu: neopravený software, slabé heslo, chybějící MFA, špatná konfigurace, neotestované zálohy nebo neproškolení zaměstnanci.

Hrozby lze stručně popsat jako aktivní, pasivní, pokročilé a trvalé. Aktivní hrozba přímo škodí, například malware nebo DDoS. Pasivní hrozba sleduje nebo sbírá informace, například odposlech provozu. Pokročilá hrozba používá sofistikované metody, třeba zero-day nebo cílený spear phishing. Trvalá hrozba odpovídá APT logice, tedy dlouhodobému a skrytému působení.

Riziko je možnost, že hrozba využije zranitelnost a způsobí škodu konkrétnímu objektu. Zjednodušeně se hodnotí podle pravděpodobnosti a dopadu, ale vždy se bere v úvahu hodnota aktiva, hrozba, zranitelnosti a připravenost organizace.

Řetězec je nejlepší vysvětlit na nemocnici. \textbf{Aktivum} je nemocniční informační systém. \textbf{Hrozba} je ransomwarová skupina. \textbf{Zranitelnost} je vzdálený přístup bez MFA. \textbf{Škodlivá událost} je průnik a zašifrování systémů. \textbf{Dopad} je omezení zdravotní péče a ohrožení pacientů. \textbf{Riziko} je možnost, že tento scénář nastane a způsobí závažné následky.

\pozor{Hrozba je možný zdroj škody, zranitelnost je slabina, riziko je možnost a závažnost škody a incident je událost, která už nastala.}

\subsection*{Vnitřní a vnější bezpečnost}
Vnitřní bezpečnost znamená ochranu pořádku, obyvatelstva a fungování státu před kriminalitou, terorismem, extremismem, korupcí a dalšími hrozbami působícími uvnitř státu. V kyberprostoru sem patří například kyberkriminalita, podvody, útoky na nemocnice, úniky dat nebo narušení digitálních služeb veřejné správy.

Vnější bezpečnost znamená ochranu před zahraniční vojenskou agresí, špionáží, nátlakem cizích států a státem podporovanými kybernetickými operacemi. V kyberprostoru sem patří například kyberšpionáž cizího státu, příprava útoků proti kritické infrastruktuře nebo kybernetické operace doprovázející ozbrojený konflikt.

V kyberprostoru se toto rozdělení stírá. Útok může být veden ze zahraničí, infrastruktura může být v několika státech, dopad se projeví uvnitř České republiky a incident současně řeší policie, NÚKIB, zpravodajské služby, diplomacie a někdy armáda.

\priklad{APT31 proti Ministerstvu zahraničních věcí ČR: česká vláda 28. 5. 2025 veřejně atribuovala dlouhodobou škodlivou kybernetickou kampaň proti jedné z neutajovaných sítí MZV ČR Čínské lidové republice; podle NÚKIB za ní velmi pravděpodobně stála skupina APT31 spojovaná s čínskou zpravodajskou službou. Příklad ukazuje prolínání vnější bezpečnosti, protože šlo o cizí státní moc, a vnitřní bezpečnosti, protože napadena byla česká veřejná instituce a dopady řešily české orgány.}

\subsection*{Bezpečnostní sektory}
Kodaňská škola rozlišuje pět bezpečnostních sektorů:
\begin{itemize}
  \item \textbf{Vojenský sektor} řeší ozbrojené síly, obranu a schopnost odstrašení; kyberpříkladem je útok na vojenské komunikační sítě.
  \item \textbf{Politický sektor} řeší suverenitu, legitimitu a fungování institucí; kyberpříkladem je útok na ministerstvo, volby nebo systémy veřejné správy.
  \item \textbf{Ekonomický sektor} řeší finance, výrobu, obchod a dodavatelské řetězce; kyberpříkladem je ransomware proti bance, továrně nebo energetické firmě.
  \item \textbf{Societální sektor} řeší identitu, sociální soudržnost a důvěru; kyberpříkladem jsou dezinformační kampaně nebo úniky dat, které podkopávají důvěru ve stát.
  \item \textbf{Environmentální sektor} řeší životní prostředí a přírodní zdroje; kyberpříkladem je útok na monitorovací systémy vody, přehrad nebo průmyslových provozů s ekologickým dopadem.
\end{itemize}

Zdravotnictví, energetika, doprava nebo komunikace nejsou dalšími teoretickými sektory Kodaňské školy. Jsou to praktická odvětví, ve kterých se několik sektorů překrývá. Útok na energetiku může být ekonomickou hrozbou kvůli výpadku výroby, politickou hrozbou kvůli oslabení fungování státu, societální hrozbou kvůli panice a vojenskou hrozbou, pokud je součástí konfliktu.

\subsection*{Sekuritizace a bezpečnostní dilema}
Sekuritizace je proces, kdy je určité téma označeno za existenční bezpečnostní hrozbu a relevantní publikum přijme potřebu mimořádných opatření. Nestačí tedy, že někdo řekne "tohle je hrozba". Důležité je, zda publikum přijme, že jde o mimořádný problém vyžadující mimořádnou reakci.

Základní prvky jsou jednoduché: \textbf{sekuritizující aktér} téma označuje za hrozbu, \textbf{referenční objekt} je to, co má být chráněno, \textbf{hrozba} je to, před čím se chrání, \textbf{publikum} musí rámování přijmout a výsledkem jsou \textbf{mimořádná opatření}. Funkční aktéři, například média, experti, dodavatelé technologií nebo regulované firmy, mohou výsledek ovlivnit, i když nejsou hlavním mluvčím.

Stupnice vypadá takto: \textbf{nepolitizované téma} stát skoro neřeší, \textbf{politizované téma} je předmětem běžné veřejné politiky, \textbf{sekuritizované téma} je prezentováno jako existenční hrozba a \textbf{desekuritizace} znamená návrat z mimořádného bezpečnostního režimu do běžné politiky.

Přednáška zdůrazňuje dva extrémy. \textbf{Sebeuspokojení} znamená podcenění reálné hrozby, například když nemocnice roky ignoruje slabé zabezpečení. \textbf{Paranoia} znamená označování neexistujících nebo přehnaných problémů za existenční hrozby. Dobrá bezpečnostní politika má mezi těmito krajnostmi udržet rozumnou rovnováhu.

Bezpečnostní dilema vzniká, když opatření jednoho aktéra ke zvýšení vlastní bezpečnosti vyvolá obavy druhého aktéra, který reaguje dalšími opatřeními, a tím vzniká spirála nedůvěry. V kyberprostoru to může být budování ofenzivních schopností nebo hromadění informací o zranitelnostech: jeden stát tvrdí, že se jen připravuje na obranu, druhý to může číst jako přípravu útoku.

\subsection*{Bezpečnostní politika}
Bezpečnostní politika je nejobecnější a vrcholný program státu nebo uskupení států v oblasti obhajoby a obrany národních zájmů. Vychází z definice národních zájmů, analýzy bezpečnostního prostředí, klasifikace hrozeb a rizik a z mezinárodních závazků. Věnuje pozornost vnitřní i vnější bezpečnosti a deklaruje, jakým způsobem a při použití jakých prostředků chce aktér, který bezpečnostní politiku definuje a provádí, hájit své zájmy. Není to jen jeden dokument; zahrnuje strategické dokumenty, právní předpisy, instituce, konkrétní opatření, financování a mezinárodní spolupráci.

V terminologii přednášky má bezpečnostní politika ČR pět základních komponent:
\begin{itemize}
  \item zahraniční politika v oblasti bezpečnosti státu,
  \item obranná politika,
  \item politika v oblasti vnitřní bezpečnosti,
  \item hospodářská politika v oblasti bezpečnosti státu,
  \item politika veřejné informovanosti v oblasti bezpečnosti státu.
\end{itemize}

Kybernetická bezpečnost je průřezová politika, protože zasahuje do všech těchto složek. Má zahraničněpolitickou rovinu při atribuci a sankcích, obrannou rovinu u kyberobrany, vnitřněbezpečnostní rovinu u kyberkriminality, hospodářskou rovinu u ochrany firem a kritické infrastruktury a informační rovinu u varování, vzdělávání a komunikace s veřejností.

\subsection*{Analýza bezpečnostní politiky}
Analýza bezpečnostní politiky se dá dělat podle jednoduché šablony:
\begin{enumerate}
  \item \textbf{Co chráníme} - referenční objekt a zájmy.
  \item \textbf{Před čím} - hrozby, zranitelnosti a rizika.
  \item \textbf{Kdo jedná} - stát, instituce, firmy, občané a útočníci.
  \item \textbf{Jaké prostředky používá} - právo, technologie, finance, diplomacii, trestní řízení a vzdělávání.
  \item \textbf{Jaké jsou náklady a vedlejší dopady} - práva, soukromí, ekonomika a inovace.
  \item \textbf{Jak poznáme, že politika funguje} - měřitelné výsledky, odolnost a schopnost obnovy.
\end{enumerate}

Tři roviny politiky stačí vysvětlit jednoduše. \textbf{Polity} je institucionální rámec, tedy kdo má pravomoci a jaké instituce existují. \textbf{Politics} je proces rozhodování a střet zájmů. \textbf{Policy} je konkrétní obsah opatření, například zákon, strategie, povinnost hlášení incidentů nebo bezpečnostní standard.

Příklad analýzy kyberbezpečnostní politiky nemocnice: chráníme dostupnost zdravotní péče, pacientská data a důvěru pacientů. Hrozbou je ransomware, phishing, zneužití dodavatele nebo výpadek systému. Zranitelností může být vzdálený přístup bez MFA, starý systém, slabá segmentace nebo neotestované zálohy. Aktéři jsou nemocnice, zřizovatel, dodavatelé, NÚKIB, policie, ÚOOÚ, ministerstvo zdravotnictví a pacienti. Nástroje jsou bezpečnostní opatření, školení, segmentace, zálohy, hlášení incidentů, trestní řízení, financování a krizová komunikace. Vedlejší dopady jsou náklady, omezení provozu a zásahy do soukromí. Fungování politiky poznáme podle toho, zda nemocnice incident včas zjistí, omezí, nahlásí, obnoví péči a sníží dopad na pacienty.

\zkouska{Bezpečnostní studia zkoumají, co je ohroženo, kdo nebo co to ohrožuje a jak lze bezpečnost chránit nebo obnovit. Vznikla jako podobor mezinárodních vztahů zaměřený hlavně na válku a vojenskou bezpečnost, ale po studené válce se rozšířila na ekonomické, politické, společenské, ekologické, kriminální a kybernetické hrozby a po 11. září 2001 i na vnitřní bezpečnost. Bezpečnost chápu jako stav, kdy jsou hrozby pro určitý objekt omezeny na přijatelnou úroveň a objekt je připraven zbylé hrozby zvládat. V kyberbezpečnosti se chrání systémy, sítě, aplikace, data a uživatelé pomocí právních, organizačních, technických a vzdělávacích opatření; prakticky jde hlavně o důvěrnost, integritu, dostupnost, odolnost a obnovu. Hrozba je možný zdroj škody, zranitelnost je slabina, riziko je možnost, že hrozba zranitelnost využije a způsobí škodu, a incident je událost, která už nastala. Vnitřní bezpečnost chrání pořádek a fungování státu uvnitř, vnější bezpečnost chrání před zahraniční agresí, špionáží a nátlakem; v kyberprostoru se ale prolínají, protože útok může přijít ze zahraničí, dopadnout uvnitř ČR a řešit ho policie, NÚKIB, zpravodajské služby i diplomacie. Kodaňská škola rozlišuje vojenský, politický, ekonomický, societální a environmentální sektor, ale praktická odvětví jako zdravotnictví nebo energetika tyto sektory kombinují. Sekuritizace znamená označení tématu za existenční hrozbu a přijetí mimořádných opatření publikem. Bezpečnostní politika je nejobecnější a vrcholný program státu nebo uskupení států pro obhajobu a obranu národních zájmů; vychází z národních zájmů, analýzy bezpečnostního prostředí, hrozeb, rizik a mezinárodních závazků a určuje prostředky pro vnitřní i vnější bezpečnost. U kyberbezpečnosti ji analyzuji přes referenční objekt, hrozby, aktéry, nástroje, náklady a ověření, zda opatření skutečně zvyšují odolnost a schopnost obnovy.}

\zdroje{Přednáškový materiál \emph{1.) Vymezení bezpečnostních studií (LONG)}; přiložené PDF v4.2; NÚKIB k veřejné atribuci APT31 proti MZV ČR; NIS2 definice rizika a bezpečnosti sítí a informačních systémů; Buzan, Waever, de Wilde: \emph{Security: A New Framework for Analysis}.}

\newpage
